% Five parts: the question; what was known; what we did; why the reader should care; why it matters.
\begin{abstract}
Do an Apple Pencil on an iPad and a pen on a Wacom tablet measure the same drawing movement?
Drawing and handwriting features are studied as health measurements, and models and reference values from one device are often applied to recordings from the other.
Earlier comparisons used people, who move differently on each surface, so the difference between the measurement chains was not known.
We built a bench in which a motion system draws the same programmed loops, spirals, and oscillating lines on both devices, while a camera at 240 frames per second measures the movement that the stylus makes.
For six features, we tested equivalence against limits that were fixed before the main runs.
Where published data allowed it, the limit is the device bias that alone gives a screening model an area under the curve of 0.55 when cases and controls use different devices.
The design has six mountings per device and 720 main runs. A cross-device correction and a simulation with generated cohorts complete the study.
\pending{State the results: which features are equivalent, the largest bias with its interval, the correction check, and the primary consequence summaries. State the tested range.}
A user of drawing features can read from these results which features transfer between the two devices, which need a correction, and how large the error of a screening model can be when devices are mixed.
The method applies to any pair of pen digitisers and gives device comparisons a mechanical reference that studies with people cannot provide.
\end{abstract}
